
\section*{Introduction}
In this chapter of the appendix we will cover some of the formalities regarding set-theory and cardinal arithmetic. The reader should notice that all results covered in this chapter (in particular the construction \ref{construction: Beth-numbers and existence of strong limit cardinals}) are well-placed within $\categoryname{ZFC}$. In particular, the approach chosen by Clausen and Scholze to approximate sheaves on a large site by approximating this large site from below using \quote{cutoff-cardinals} is well-placed too and provides a good alternative to usual approaches like working with Grothendieck universes -- which usually require additional axioms to be assumed on top of $\categoryname{ZFC}$. We are lenient in the uses of \emph{classes} as these are only used to group objects for ease of comprehension, and no result stated fundamentally depends on the existence of classes as objects on their own right.

\begin{convention}
	We will use the \emph{Von Neumann cardinal assignment}. That is, we choose for each class of equicardinal sets a specific representative, namely the smallest (recall the class of all ordinals is well-ordered) ordinal in this equivalence class. Denoting by $\Ordinals$ the class of all ordinals, we thus define for any set $S$
		$$\vert S\vert \ldef \inf\{\alpha \in \Ordinals \setseparator \alpha \sim S\}$$
	where $S\sim S'$ if and only if there exists a bijection between $S$ and $S'$. We define the class of cardinal numbers to be the class of these representatives. Hence,
		$$\Cardinals \ldef \{\vert S\vert \setseparator S \in \Set \}.$$
	In particular every cardinal is an ordinal and the class $\Cardinals$ of cardinals is well-ordered itself.
\end{convention}

\begin{definition}[Cardinal arithmetic]
	For a cardinal $\kappa$ define its \emph{successor cardinal} as the cardinal
		$$\kappa^+ \ldef \inf\{\kappa' \in \Cardinals \setseparator \kappa' > \kappa\}.$$
	For a second cardinal $\kappa'$ write
		$$\kappa + \kappa' \ldef \vert \kappa \sqcup \kappa' \vert$$
	where $\sqcup$ denotes the coproduct in $\Set$, as well as
		$$\kappa \cdot \kappa' \ldef \vert \kappa \times \kappa'\vert.$$
	The cardinality of the powerset is denoted
		$$2^\kappa \ldef \vert \powerset \kappa\vert.$$
	More generally for cardinals $\mu$ and $\lambda$ define
		$$\mu^\lambda \ldef \vert \{f\colon \lambda \to \mu\}\vert$$
	to be the cardinality of the set of all functions $\lambda\to \mu$.
\end{definition}

\begin{definition}[Weak \& strong limit cardinals]
	A cardinal $\kappa$ is a \emph{weak limit cardinal} if it is neither $0$ nor a successor cardinal. If in addition for every $\kappa' < \kappa$ also $2^{\kappa'} < \kappa$ then $\kappa$ is called a \emph{strong limit cardinal}. 
\end{definition}

\begin{construction}[$\beth$-numbers \& strong limit cardinals]
	\label{construction: Beth-numbers and existence of strong limit cardinals}	
	Define recursively
		$$\beth_0 \ldef 0,$$
	for every ordinal $\alpha$
		$$\beth_{\alpha^+} \ldef 2^{\beth_\alpha}$$
	and finally for $\alpha$ a limit ordinal
		$$\beth_\alpha \ldef \bigcup_{\alpha' < \alpha} \beth_{\alpha'}.$$
	Denoting by $\omega$ the order-type of $\N$, we obtain for each ordinal $\alpha$ an uncountable large strong limit cardinal
		$$\beth_{\omega + \alpha} = \beth_{\alpha + \omega} = \bigcup_{n < \omega} \beth_{\alpha + n},$$
	in particular we can construct arbitrarily large uncountable strong limit cardinals.
\end{construction}

\begin{definition}[Cofinality for category theorists]
	Let $\kappa$ be a cardinal. The \emph{cofinality} $\cofinality\kappa$ of $\kappa$ is the smallest cardinal $\alpha$ such that $\boundedSet \kappa$ is not $\alpha$-cocomplete. Hence for every cardinal $\lambda$ we have that $\lambda < \cofinality{\kappa}$ if and only if $\boundedSet \kappa$ is $\lambda$-cocomplete.
\end{definition}

\begin{remark}[An upper bound]
	\label{remark: upper bound of cofinality}
	
	Let $\kappa$ be a cardinal. Then $\kappa$ is a colimit of a diagram of size $\kappa$ ($\kappa$ is in bijection with the disjoint union of its elements.) Thus, as $\kappa$ is not $\kappa$-small, the category $\boundedSet \kappa$ is not $\kappa$-cocomplete. Hence, $\cofinality{\kappa} \le \kappa$.
\end{remark}

\begin{lemma}
	For any infinite cardinal $\kappa$ its cofinality $\cofinality{\kappa}$ is infinite as well.
\end{lemma}
\begin{proof}
	Suppose for a contradiction that $\cofinality{\kappa} = n \in \N$ is finite. Then by definition $\boundedSet{\kappa}$ is not $n$-complete. Hence, there are $n$ cardinals $\kappa_1,\dots,\kappa_n$ smaller than $\kappa$ with
		$$\sup\{\kappa_1,\cdots,\kappa_n\} \ge \kappa.$$
	But the supremum of finitely many cardinals is just their maximum. Thus,
		$$\sup\{\kappa_1,\dots,\kappa_n\} < \kappa$$
	in contradiction to the previous inequality.
\end{proof}

\begin{definition}[Regular and singular cardinals]
	A cardinal $\kappa$ is called \emph{regular} if $\cofinality{\kappa} = \kappa$. Hence, $\kappa$ is regular if and only if the category $\boundedSet \kappa$ is $\lambda$-complete for all $\lambda < \kappa$. If $\kappa$ is not regular it is called \emph{singular}.
\end{definition}

\begin{example}[Regular and singular cardinals]
	Certainly the category $\set$ of finite sets is closed under finite colimits. Hence, $\N$ is regular. To construct a singular cardinal is more work. Construct the $\aleph$-sequence by defining recursively
		$$\aleph_0\ldef\N$$
	then
		$$\aleph_{\alpha^+} \ldef \inf\{\kappa \in \Cardinals\setseparator\kappa > \aleph_\alpha\}$$
	for every ordinal $\alpha$, recalling that cardinals are well-ordered and finally
		$$\aleph_\alpha \ldef \bigcup_{\lambda < \alpha}\aleph_\lambda$$
	for a limit ordinal $\alpha$. Then $\aleph_\omega$ is the union of the sequence $(\aleph_n)_{n\in\N}$ (which is the colimit of the diagram $n \to \aleph_n$ with $\N$ as a poset category). Hence, $\boundedSet{\aleph_\omega}$ is not even $\N$-cocomplete. In particular $\cofinality{\aleph_\omega}\le \N \ll \aleph_\omega$ and thus $\aleph_\omega$ is singular.
\end{example}

%\begin{lemma}[Cofinalities are regular]
%	\label{lemma: cofinalities are regular}
%	For any cardinal $\kappa$ the cofinality $\cofinality{\kappa}$ is a regular cardinal.
%\end{lemma}
%\begin{proof}
%	By remark \ref{remark: upper bound of cofinality} certainly $\cofinality{\cofinality{\kappa}}\le \cofinality{\kappa}$. Assume that $\cofinality{\cofinality{\kappa}} < \cofinality{\kappa}$.
%	\begin{todo}
%		% TODO:
%		Proof or reference, but keeping in mind that this definition of cofinality is a bit non-standard.
%	\end{todo}
%\end{proof}
